% guia_google_analytics.tex
% Guía adicional: Google Analytics 4 (contexto, usos, visualización y extracción de datos)
% Laboratorio: Gestión de datos, comportamiento y perfiles del consumidor digital
% Lic. en Mercadotecnia Digital · ESCA Unidad Santo Tomás · IPN
% Fuentes: ayuda oficial de Google Analytics y documentación para desarrolladores (consultadas el 12-oct-2026)
\documentclass[11pt,letterpaper]{article}
\usepackage[utf8]{inputenc}
\usepackage[T1]{fontenc}
\usepackage{lmodern}
\usepackage[spanish,es-noshorthands,es-nodecimaldot]{babel}
\usepackage[margin=2.4cm,headheight=14pt]{geometry}
\usepackage{csquotes}
\usepackage{xcolor}
\usepackage{graphicx}
\usepackage{booktabs}
\usepackage{array}
\usepackage{enumitem}
\usepackage{fancyhdr}
\usepackage{xurl}
\usepackage[hidelinks]{hyperref}
\usepackage[most]{tcolorbox}
\usepackage{listings}
\usepackage{tabularx}
\usepackage{caption}
\captionsetup{font=small,labelfont=bf}

\definecolor{guinda}{HTML}{6F1D46}
\definecolor{guindaOsc}{HTML}{45102C}
\definecolor{guindaClara}{HTML}{F4E9EF}
\definecolor{dorado}{HTML}{B8975A}
\definecolor{terracota}{HTML}{9C5B44}
\definecolor{bronce}{HTML}{8C6B3F}
\definecolor{gris}{HTML}{58595B}
\hypersetup{colorlinks=true,linkcolor=guinda,urlcolor=guinda,citecolor=guinda,
  pdftitle={Guia: Google Analytics 4 para mercadotecnia digital},pdfauthor={Dr. Aboud Barsekh-Onji}}

\newtcolorbox{cajaNota}[1][]{colback=guindaClara,colframe=guinda,breakable,fonttitle=\bfseries,title={Idea clave},#1}
\newtcolorbox{cajaAtencion}[1][]{colback=terracota!8,colframe=terracota,breakable,fonttitle=\bfseries,title={Atención},#1}
\newtcolorbox{cajaPasos}[1][]{colback=bronce!7,colframe=bronce,breakable,fonttitle=\bfseries,title={Paso a paso},#1}
\newtcolorbox{cajaVigencia}[1][]{colback=gris!7,colframe=gris,breakable,fonttitle=\bfseries,title={Vigencia},#1}

\lstdefinestyle{csv}{basicstyle=\ttfamily\scriptsize,breaklines=true,frame=single,backgroundcolor=\color{gray!6},
  rulecolor=\color{guinda!60},columns=fullflexible,keepspaces=true}
\lstset{style=csv,literate={á}{{\'a}}1 {é}{{\'e}}1 {í}{{\'i}}1 {ó}{{\'o}}1 {ú}{{\'u}}1 {ñ}{{\~n}}1}

\pagestyle{fancy}\fancyhf{}
\fancyhead[L]{\small\color{gris}Google Analytics 4 para mercadotecnia digital}
\fancyhead[R]{\small\color{gris}IPN · ESCA Santo Tomás}
\fancyfoot[C]{\small\thepage}
\renewcommand{\headrulewidth}{0.3pt}
\setlength{\parindent}{0pt}\setlength{\parskip}{5pt}
\renewcommand{\arraystretch}{1.15}

\title{\textbf{Google Analytics 4 para mercadotecnia digital}\\[4pt]
  \large Contexto, usos y pasos para visualizar y sacar datos}
\author{Dr. Aboud Barsekh-Onji\\[2pt]\small Instituto Politécnico Nacional · ESCA Unidad Santo Tomás\\
  \small Laboratorio: Gestión de datos, comportamiento y perfiles del consumidor digital}
\date{Octubre de 2026}

\begin{document}

\begin{center}
  \includegraphics[height=1.6cm]{ipn-escudo.png}\\[6pt]
  {\Large\bfseries\color{guinda} Google Analytics 4 para mercadotecnia digital}\\[3pt]
  {\large Contexto, usos y pasos para visualizar y sacar datos}\\[8pt]
  {\small Dr. Aboud Barsekh-Onji · Instituto Politécnico Nacional · ESCA Unidad Santo Tomás\\
  Laboratorio: Gestión de datos, comportamiento y perfiles del consumidor digital\\
  Material adicional de la Unidad I · octubre de 2026}
\end{center}
\vspace{2pt}\hrule\vspace{6pt}

\begin{cajaVigencia}
Este documento se redactó con la ayuda oficial de Google Analytics y la documentación para desarrolladores \textbf{consultadas el 12 de octubre de 2026}.
Tal como advierte el Capítulo~4, \textbf{las interfaces, los nombres de los menús y los límites de una plataforma cambian más rápido que cualquier texto}: antes de
aplicar un paso en una propiedad real, compáralo con la página de ayuda vigente (se citan al final). Los nombres de menús se dan en inglés, como aparecen en la
documentación; la interfaz en español puede traducirlos de otra forma.
\end{cajaVigencia}

\tableofcontents
\newpage

%=====================================================================
\section{Para qué sirve esta guía}

En la Unidad~I se estudió quién es el consumidor digital, cómo se perfila, cómo se gestionan sus datos y qué criterios sirven para evaluar plataformas. \textbf{Google Analytics}
es una de las herramientas con las que una organización \emph{observa} ese comportamiento en su propio sitio web o app: ejemplifica la \enquote{señal conductual} y los
\enquote{datos de primera parte} de la Sesión~3 y produce las series de tiempo que se analizaron en la Sesión~4.

Al terminar esta guía podrás:
\begin{enumerate}[nosep]
  \item explicar qué es Google Analytics 4 (GA4), de dónde viene y cómo recolecta datos;
  \item relacionar preguntas de mercadotecnia con los informes de GA4 que las responden;
  \item \textbf{visualizar} datos en los informes estándar y en las Exploraciones, usando la cuenta de demostración;
  \item \textbf{sacar} datos de GA4 (CSV y Looker Studio; BigQuery y la API solo como referencia), sabiendo que la cuenta de demostración no permite exportar y que sí se podrá en una cuenta propia;
  \item identificar las obligaciones de privacidad que implica todo lo anterior.
\end{enumerate}

%=====================================================================
\section{Contexto: qué es Google Analytics y de dónde viene}

Google Analytics es un servicio gratuito de analítica web y de aplicaciones: registra las interacciones de las personas con un sitio o una app y las resume en informes,
con una versión de pago (Analytics 360) para organizaciones grandes. Sus dos generaciones relevantes son:

\begin{table}[h]
\centering
\begin{tabular}{@{}p{3.3cm}p{5.6cm}p{5.6cm}@{}}
\toprule
 & \textbf{Universal Analytics (UA)} & \textbf{Google Analytics 4 (GA4)} \\
\midrule
Unidad de medida & Sesión y \enquote{vistas de página} & \textbf{Evento} (con parámetros) \\
Estado & Retirada & Vigente; reemplazó a UA \\
Plataformas & Sitios web (apps con otra herramienta) & Sitios web y apps en una misma propiedad \\
\bottomrule
\end{tabular}
\caption*{Fuente: elaboración propia con base en la ayuda de Google Analytics (ver referencias).}
\end{table}

\begin{cajaNota}[title={Cronología de la transición (fuente: ayuda oficial de Google)}]
\begin{itemize}[nosep]
  \item \textbf{1 de julio de 2023:} las propiedades estándar de Universal Analytics dejaron de procesar datos nuevos.
  \item \textbf{1 de julio de 2024:} se perdió el acceso a la interfaz y a la API de UA, y los datos históricos de UA dejaron de estar disponibles.
  \item \textbf{Hoy:} GA4 es la solución de medición vigente. Quien hoy lea \enquote{Analytics} en un tutorial viejo puede estar viendo UA: conviene revisar la fecha.
\end{itemize}
\end{cajaNota}

\textbf{Por qué importa para la mercadotecnia digital.} El cambio de \emph{sesiones} a \emph{eventos} no es solo técnico: obliga a decidir \emph{qué comportamiento se mide}
(¿un clic?, ¿ver un video?, ¿completar una compra?). Es la misma idea de la Sesión~2: el perfil depende de qué comportamiento se observa.

%=====================================================================
\section{Cómo funciona GA4: el modelo de datos}
\label{sec:modelo}

\subsection{Eventos, parámetros y propiedades de usuario}

En GA4 todo es un \textbf{evento}: una interacción con un nombre (por ejemplo, \texttt{page\_view} o \texttt{purchase}) y un conjunto de \textbf{parámetros} que la describen
(la página, el valor de la compra, la moneda\ldots). Además se pueden asociar \textbf{propiedades de usuario} (atributos de la persona, como un segmento) y, en el comercio
electrónico, \textbf{parámetros de artículo} (producto, categoría, precio).

\begin{table}[h]
\centering
\begin{tabular}{@{}p{3.6cm}p{6.2cm}p{4.6cm}@{}}
\toprule
\textbf{Categoría de evento} & \textbf{Qué es} & \textbf{Ejemplos} \\
\midrule
Automáticos & Se registran al instalar la etiqueta o el SDK, sin configurar nada & \texttt{first\_visit}, \texttt{session\_start}, \texttt{user\_engagement}, \texttt{first\_open} (app) \\
Medición mejorada & Se activan por defecto en flujos web & \texttt{page\_view}, \texttt{scroll}, \texttt{click}, \texttt{file\_download}, \texttt{form\_start}, \texttt{form\_submit}, \texttt{video\_start} \\
Recomendados & Nombres sugeridos por Google, que hay que implementar & \texttt{purchase}, \texttt{add\_to\_cart}, \texttt{login} (ver ayuda de eventos recomendados) \\
Personalizados & Los define la organización & Cualquier nombre propio, p.\,ej. \texttt{compra\_en\_tienda} \\
\bottomrule
\end{tabular}
\caption*{Fuente: elaboración propia con base en la ayuda \enquote{Automatically collected events} y la de eventos recomendados.}
\end{table}

\subsection{Dimensiones y métricas}

Los informes combinan dos tipos de campos. Una \textbf{dimensión} es un \emph{atributo} que sirve para agrupar (país, canal de adquisición, página, dispositivo).
Una \textbf{métrica} es una \emph{cantidad} que se mide (usuarios activos, sesiones, número de eventos, eventos clave, ingresos). Un informe es, en esencia, una tabla:
una columna por dimensión y una por métrica. Las dimensiones personalizadas pueden tener ámbito de \textbf{evento}, de \textbf{usuario} o de \textbf{artículo}, y deben crearse
\emph{antes} de poder usarlas.

\subsection{Cómo se recolectan los datos}

\begin{itemize}[nosep]
  \item \textbf{Etiqueta de Google (\emph{Google tag}) en sitios web}, instalada directamente o mediante Google Tag Manager.
  \item \textbf{SDK de Firebase en apps} móviles.
  \item \textbf{Measurement Protocol}: envío de eventos desde un servidor u otro sistema que no puede usar la etiqueta (p.\,ej., ventas fuera de línea).
\end{itemize}
Cada fuente se declara como un \textbf{flujo de datos} (\emph{data stream}) de la propiedad, con su propio identificador de medición.

\subsection{Cómo se identifica a una persona (sin saber quién es)}

GA4 no \enquote{conoce} a nadie: reconoce \emph{dispositivos y navegadores} mediante un \textbf{identificador de cliente} (\texttt{client\_id}) guardado en la cookie \texttt{\_ga}, cuya
caducidad por defecto es de \textbf{2 años} (configurable de inmediato hasta 25 meses; los navegadores pueden acortarla). Si la organización tiene inicios de sesión, puede enviar
además un \textbf{User-ID} propio. Una persona con tres dispositivos y sin sesión iniciada aparece como tres usuarios: es el problema de la identidad única de la omnicanalidad (Sesión~3).

\subsection{Límites de recolección que conviene conocer}

Para propiedades estándar, la ayuda de Google indica, entre otros: hasta \textbf{25 parámetros por evento}, \textbf{25 propiedades de usuario}, nombres de evento y de parámetro de
\textbf{40 caracteres como máximo} (alfanuméricos y guion bajo, empezando con letra), valores de parámetro de \textbf{100 caracteres}, valores de propiedad de usuario de \textbf{36 caracteres} y
User-ID de hasta 256 caracteres. Son los límites de lo que se recolecta.

%=====================================================================
\section{Para qué se usa en mercadotecnia digital}

\begin{table}[h]
\centering
\small
\begin{tabularx}{\textwidth}{@{}>{\raggedright\arraybackslash}p{4.4cm}>{\raggedright\arraybackslash}X>{\raggedright\arraybackslash}p{2.9cm}@{}}
\toprule
\textbf{Pregunta de negocio} & \textbf{Dónde se responde en GA4} & \textbf{Se conecta con} \\
\midrule
¿De dónde llegan mis visitantes y cuáles canales convierten? & Informes de \emph{Acquisition} (adquisición de usuarios y de tráfico) & Cap.~4 (plataformas, atribución) \\
¿Qué hacen en mi sitio y en qué paso abandonan? & \emph{Engagement}; Exploración de embudo (\emph{funnel}) y de rutas (\emph{path}) & Cap.~1 (reducción de fricción) \\
¿Qué se vende y cuánto ingreso genera? & \emph{Monetization}; datos de artículos & Cap.~3 (valor) \\
¿Vuelven los usuarios? ¿Cuánto vale una cohorte? & \emph{Retention}; Exploración de cohortes y de valor de vida (\emph{user lifetime}) & Cap.~3 (CLV) \\
¿Qué segmentos de usuarios tengo? & Audiencias; Exploración de superposición de segmentos & Cap.~2 (segmentación) \\
¿Hay estacionalidad en mi tráfico? & Series diarias o semanales leídas y descompuestas & Cap.~4 (laboratorio) \\
\bottomrule
\end{tabularx}
\caption*{Fuente: elaboración propia.}
\end{table}

\begin{cajaNota}
\textbf{Un dato puede ser tan bueno como la pregunta.} GA4 mide comportamiento \emph{agregado y anónimo} dentro de \emph{su} sitio o app: no dice por qué la gente actúa así
(señal declarada) ni qué hace fuera de ese sitio. Complementa --- no sustituye --- encuestas, escucha social y datos propios (CRM).
\end{cajaNota}

%=====================================================================
\section{Cómo visualizar los datos}
\label{sec:visualizar}

\subsection{Acceso para practicar: la cuenta de demostración}

Google ofrece una \textbf{cuenta de demostración} con datos reales de comercio electrónico de la tienda de mercancía de Google (\emph{Google Merchandise Store}) y de un juego (\emph{Flood-It!}).
Es de solo lectura y no requiere instalar nada.

\begin{cajaPasos}
\begin{enumerate}[nosep,label=\arabic*.]
  \item Inicia sesión con una cuenta de Google y abre el enlace oficial de la demostración: \url{https://goo.gle/ga-demo}.
  \item Acepta los términos. Si no tenías cuenta de Analytics, Google crea una y le agrega la cuenta de demostración.
  \item En el selector de cuentas/propiedades (arriba a la izquierda), elige la propiedad \textbf{GA4 -- Google Merchandise Store}.
  \item Verifica en \emph{Admin} que ves la propiedad como \emph{Viewer} (lector): no podrás cambiar la configuración, y eso es lo esperado.
\end{enumerate}
\end{cajaPasos}

\begin{cajaAtencion}
La cuenta de demostración cuenta contra el máximo de cuentas de Analytics que puedes tener con un mismo usuario de Google, y puedes quitarla cuando quieras. Verifica en el enlace oficial
que las propiedades siguen disponibles: son datos de Google, no del curso.
\end{cajaAtencion}

\subsection{Mapa de la interfaz}

\begin{table}[h]
\centering
\begin{tabular}{@{}p{3.2cm}p{11.3cm}@{}}
\toprule
\textbf{Sección} & \textbf{Para qué sirve} \\
\midrule
Home & Resumen y tarjetas con lo más reciente \\
Reports & Informes estándar: \emph{Reports snapshot}, \emph{Realtime} y colecciones de informes \\
Explore & Exploraciones: análisis libres y avanzados que tú configuras \\
Advertising & Atribución y planeación de canales \\
Admin & Configuración: flujos de datos, eventos clave, dimensiones personalizadas, retención, vínculos \\
\bottomrule
\end{tabular}
\caption*{Fuente: elaboración propia con base en la ayuda \enquote{Reports in the Analytics app}.}
\end{table}

\subsection{Informes estándar (\emph{Reports})}

Los informes están organizados en \textbf{colecciones}. La colección predeterminada \emph{Life cycle} sigue el recorrido del cliente (\emph{Acquisition}, \emph{Engagement},
\emph{Monetization}, \emph{Retention}) y la colección \emph{User} resume \emph{User attributes} y \emph{Tech}. El \emph{Reports snapshot} es el resumen general y
\emph{Realtime} muestra la actividad de los últimos minutos. Cada tema tiene un informe de resumen (\emph{overview}) y informes de detalle.

\begin{cajaPasos}[title={Práctica: ¿de dónde llega el tráfico?}]
\begin{enumerate}[nosep,label=\arabic*.]
  \item En \emph{Reports}, abre \textbf{Acquisition $\rightarrow$ Traffic acquisition}.
  \item Fija el rango de fechas (arriba a la derecha) a los últimos 28 días.
  \item Cambia la dimensión principal a \emph{Session source / medium} con el selector de la primera columna.
  \item Agrega una comparación (\emph{Add comparison}) para ver un segmento, por ejemplo un tipo de dispositivo.
  \item Identifica el canal con más sesiones y el canal con mejor tasa de interacción. ¿Coinciden?
\end{enumerate}
\end{cajaPasos}

\textbf{Eventos clave.} Las acciones que importan al negocio (compra, registro, envío de formulario) se marcan como \textbf{eventos clave} (\emph{key events}); es el
término que Google usa en Analytics, mientras que \enquote{conversión} se reserva para el contexto de anuncios. Un evento clave se marca en \emph{Admin} y luego puede
usarse en informes y exportarse a Google Ads.

\subsection{Exploraciones (\emph{Explore})}

Las Exploraciones permiten construir análisis a la medida, con \textbf{siete técnicas}:

\begin{table}[h]
\centering
\begin{tabular}{@{}p{3.6cm}p{10.9cm}@{}}
\toprule
\textbf{Técnica} & \textbf{Qué responde} \\
\midrule
Free form & Tablas cruzadas, con barras, líneas, dispersión o mapas \\
Funnel exploration & Cuántos usuarios avanzan en cada paso de una tarea y dónde se pierden \\
Path exploration & Qué rutas siguen los usuarios por el sitio o la app \\
Segment overlap & Cómo se traslapan distintos segmentos de usuarios \\
User explorer & Quiénes son los usuarios de un segmento y qué hizo cada uno \\
Cohort exploration & Cómo se comportan grupos de usuarios que comparten un atributo (p.\,ej., fecha de llegada) \\
User lifetime & Comportamiento y valor de los usuarios a lo largo de su vida como clientes \\
\bottomrule
\end{tabular}
\caption*{Fuente: elaboración propia con base en la ayuda \enquote{Explorations}.}
\end{table}

\begin{cajaPasos}[title={Práctica: un embudo de compra}]
\begin{enumerate}[nosep,label=\arabic*.]
  \item En \emph{Explore}, crea una exploración en blanco y elige la técnica \textbf{Funnel exploration}.
  \item Define los pasos con eventos: \texttt{view\_item} $\rightarrow$ \texttt{add\_to\_cart} $\rightarrow$ \texttt{begin\_checkout} $\rightarrow$ \texttt{purchase}.
  \item Anota el porcentaje de usuarios que pasa de un paso al siguiente y señala el paso con mayor abandono (la \emph{fricción} del Capítulo~1).
  \item Compara dos segmentos (por ejemplo, móvil y escritorio): ¿abandonan en el mismo paso?
\end{enumerate}
\end{cajaPasos}

\subsection{Cuánto tiempo se conservan los datos}

La retención de datos de usuario y de evento en propiedades estándar puede configurarse en \textbf{2 meses o 14 meses}. Este ajuste afecta a las \textbf{Exploraciones y a los embudos},
\emph{no} a los informes estándar agregados. Por eso una exploración de más de 14 meses \enquote{se corta} aunque el informe estándar sí muestre ese periodo. Los datos de edad, género e
intereses se conservan solo 2 meses.

\begin{cajaAtencion}
\textbf{Lectura crítica de series.} Antes de interpretar un cambio, compara con el mismo periodo del año anterior y revisa si la serie cubre al menos un ciclo anual completo:
es exactamente la pregunta del Capítulo~4 (\enquote{¿campaña o temporada?}). Un informe de 28 días no basta para distinguir una tendencia de una estacionalidad.
\end{cajaAtencion}

%=====================================================================
\section{Cómo sacar datos de GA4 y llevarlos a otras herramientas}
\label{sec:sacar}

\begin{table}[h]
\centering
\small
\begin{tabularx}{\textwidth}{@{}>{\raggedright\arraybackslash}p{2.8cm}>{\raggedright\arraybackslash}X>{\raggedright\arraybackslash}p{3.8cm}@{}}
\toprule
\textbf{Vía} & \textbf{Qué ofrece} & \textbf{Cuándo usarla} \\
\midrule
Exportar desde una exploración & Google Sheets, CSV, TSV o PDF desde \emph{Export data} (CSV, TSV y Sheets incluyen todos los datos de la visualización; los informes estándar tienen opciones de descarga o de compartir equivalentes: revisa la ayuda vigente) & Una tabla puntual para Excel o para la app del Capítulo~4 \\
Looker Studio & Tableros interactivos conectados a la propiedad & Reportes periódicos para otras personas \\
Exportación a BigQuery & Datos de evento \emph{crudos}, uno por fila, en un almacén de Google Cloud & Solo se menciona como referencia: lo configura el área de TI y puede generar costos \\
API de datos (\emph{Data API}) & Conexión automática de GA4 con otros programas & Solo como referencia: requiere programar y no se usa en este curso \\
\bottomrule
\end{tabularx}
\caption*{Fuente: elaboración propia con base en la ayuda de Google y la documentación de la API de datos.}
\end{table}

\begin{cajaNota}
\textbf{La cuenta de demostración no permite exportar.} La cuenta de demostración de Google es de \emph{solo consulta}: no se pueden descargar ni exportar datos (la opción \emph{Export data} no está disponible). Es una restricción de esa cuenta pública, no un error del usuario. En una \textbf{propiedad propia}, con un rol de Analista o superior, sí se podrá exportar a CSV, Sheets o PDF (verifica en la ayuda vigente el rol que exige cada exportación). Por eso las prácticas de esta guía \textbf{no requieren descargar nada}: se resuelven leyendo y capturando la pantalla.
\end{cajaNota}

\begin{cajaPasos}[title={Práctica: leer una serie en pantalla (tendencia y estacionalidad)}]
\begin{enumerate}[nosep,label=\arabic*.]
  \item En \emph{Explore}, crea una exploración \emph{Free form} con visualización de \textbf{líneas}: pon como dimensión la \textbf{semana} (o la fecha) y como métrica las \textbf{sesiones}; usa el rango más largo disponible (idealmente $\geq$ 2 años, y recuerda que la retención de 2 o 14 meses limita las exploraciones).
  \item Toma una \textbf{captura de pantalla} de la gráfica.
  \item Sobre la captura, marca con lápiz la \textbf{tendencia} (¿sube, baja o se mantiene?) y los \textbf{picos que se repiten} cada año.
  \item Interpreta en una o dos frases: ¿es un patrón de temporada o un salto de campaña? ¿hay dos ciclos completos para confirmarlo?
\end{enumerate}
\end{cajaPasos}

\textbf{Opcional, solo con cuenta propia.} Exporta la serie a CSV, déjala en Excel con dos columnas (fecha y valor), súbela a \texttt{cap04\_series\_datos\_propios.html} y ajusta \emph{periodos por ciclo} a 52 (semanal).

\textbf{Precaución con los datos agregados.} Lo que sale de un informe ya viene \emph{agregado} (no hay filas por persona), y los informes pueden aplicar umbrales de privacidad,
muestreo en exploraciones muy grandes y diferencias de procesamiento: dos fuentes (informe, otra exportación) pueden no coincidir exactamente. Documenta siempre de qué
fuente y con qué fecha de consulta salió cada cifra.

%=====================================================================
\section{Privacidad, ética y base legal}

\begin{itemize}[nosep]
  \item \textbf{Es tratamiento de datos personales.} Una cookie que identifica un navegador, un User-ID o un CRM unido a la analítica son datos que el Capítulo~3 sitúa bajo la LFPDPPP. Se necesita
        \textbf{base legal} y un \textbf{propósito declarado}; verifica el texto vigente de la ley, como se advirtió en la Sesión~3.
  \item \textbf{Aviso y consentimiento.} El sitio debe informar el uso de cookies y analítica. GA4 ofrece \emph{Consent Mode}, que ajusta el comportamiento de las etiquetas según el consentimiento
        del usuario (consulta la referencia oficial para su configuración).
  \item \textbf{Minimización.} Recolecta solo lo necesario; nunca subas nombres, correos, teléfonos ni direcciones: la política de Google lo prohíbe y el principio es el mismo que el de
        los datos de cero parte (lo que el consumidor entrega con conocimiento).
  \item \textbf{IP.} La documentación de Google indica que Analytics no registra ni almacena direcciones IP (y que, para el tráfico de la UE, solo se usan para derivar ubicación aproximada y se descartan de inmediato); aun así, la organización sigue siendo responsable del resto de los datos que recolecta.
  \item \textbf{Retención.} Elige 2 o 14 meses según la necesidad real, y elimina datos a solicitud cuando corresponda.
\end{itemize}

%=====================================================================
\section{Práctica guiada (60 a 90 minutos)}

\textbf{Entregable:} un reporte breve (2 a 3 páginas) con capturas de pantalla y las respuestas a las preguntas.

\begin{enumerate}
  \item \textbf{Contexto (10 min).} Con tus palabras: ¿qué diferencia hay entre una \emph{sesión} y un \emph{evento}? ¿Por qué el cambio de UA a GA4 importa a una marca de moda que vende en línea y en tienda?
  \item \textbf{Acceso (5 min).} Abre la propiedad \emph{GA4 -- Google Merchandise Store} de la cuenta de demostración. Captura la pantalla de \emph{Reports snapshot}.
  \item \textbf{Adquisición (15 min).} En \emph{Traffic acquisition}, para los últimos 28 días: ¿qué canal aporta más sesiones? ¿Cuál tiene mejor tasa de interacción? Explica en una frase la diferencia.
  \item \textbf{Embudo (15 min).} Construye el embudo \texttt{view\_item} $\rightarrow$ \texttt{add\_to\_cart} $\rightarrow$ \texttt{begin\_checkout} $\rightarrow$ \texttt{purchase}. ¿Dónde está el mayor abandono? Relaciónalo con la reducción de fricción (Cap.~1).
  \item \textbf{Leer una serie (10 min).} En \emph{Explore}, grafica en líneas las sesiones semanales del periodo más largo disponible y captura la pantalla (la demo no permite exportar). Marca tendencia y picos repetidos. ¿Qué estacionalidad encuentras? ¿Hay dos ciclos completos para confirmarla?
  \item \textbf{Reflexión (10 min).} Con la matriz de criterios del Capítulo~4 (modelo de negocio, intermediación de datos, interoperabilidad y exposición regulatoria), evalúa a Google Analytics: ¿quién es dueño del dato? ¿qué tan fácil es llevarse los datos a otra herramienta?
\end{enumerate}

\textbf{Rúbrica sugerida (100 puntos):} contexto 10 · acceso y capturas 10 · adquisición 20 · embudo 25 · lectura de la serie y estacionalidad 25 · reflexión con la matriz 10.

%=====================================================================
\section{Autoevaluación}

\begin{enumerate}
  \item \textbf{¿Cuál es la unidad básica de medición de GA4?}\quad a) la sesión;\quad b) el evento, con sus parámetros;\quad c) la página vista.
  \item \textbf{¿Qué efecto tiene la retención de datos de 14 meses?}\quad a) borra los informes estándar anteriores;\quad b) limita las Exploraciones y los embudos, no los informes estándar;\quad c) impide exportar datos.
  \item \textbf{¿Cómo vienen los datos que se exportan de un informe de GA4?}\quad a) una fila por persona;\quad b) agregados, por lo que hay que documentar su fuente;\quad c) sin ninguna restricción de privacidad.
  \item \textbf{¿Qué debes evitar siempre al recolectar o analizar datos de usuarios?}\quad a) los datos agregados;\quad b) información personal identificable como correos o teléfonos;\quad c) los informes estándar.
\end{enumerate}
\textbf{Clave:} 1\,b · 2\,b · 3\,a · 4\,b · 5\,b.

%=====================================================================
\section{Glosario}

\begin{description}[leftmargin=1.2cm,style=nextline]
  \item[Evento] Interacción con nombre y parámetros (p.\,ej. \texttt{purchase}); unidad de medición de GA4.
  \item[Parámetro] Dato que describe un evento (valor, moneda, página).
  \item[Propiedad de usuario] Atributo de una persona asociado a sus eventos (p.\,ej. segmento).
  \item[Dimensión / métrica] Atributo para agrupar / cantidad que se mide.
  \item[Flujo de datos (\emph{data stream})] Fuente de datos de la propiedad: sitio web, app Android o iOS.
  \item[Evento clave (\emph{key event})] Evento marcado como importante para el negocio (antes llamado \enquote{conversión} en Analytics).
  \item[Exploración] Análisis a la medida (tabla libre, embudo, rutas, cohortes\ldots).
  \item[Client ID / User-ID] Identificador del navegador (cookie \texttt{\_ga}) / identificador propio de la organización para usuarios con sesión.
  \item[BigQuery] Almacén de datos en la nube de Google al que GA4 puede exportar eventos crudos.
  \item[Consent Mode] Mecanismo que ajusta lo que miden las etiquetas según el consentimiento del usuario.
\end{description}

%=====================================================================
\section{Referencias y fuentes consultadas}

Todas las páginas se consultaron el 12 de octubre de 2026; la numeración corresponde al identificador de la página de ayuda.

\begin{itemize}[nosep,leftmargin=*]\small
  \item Google. (s.\,f.). \emph{Google Analytics 4 has replaced Universal Analytics}. Analytics Help. \url{https://support.google.com/analytics/answer/9973999}
  \item Google. (s.\,f.). \emph{[GA4] Automatically collected events}. Analytics Help. \url{https://support.google.com/analytics/answer/9234069}
  \item Google. (s.\,f.). \emph{[GA4] Collection limits}. Analytics Help. \url{https://support.google.com/analytics/answer/9267744}
  \item Google. (s.\,f.). \emph{[GA4] Reports in the Analytics app}. Analytics Help. \url{https://support.google.com/analytics/answer/9924671}
  \item Google. (s.\,f.). \emph{[GA4] Explorations}. Analytics Help. \url{https://support.google.com/analytics/answer/7579450}
  \item Google. (s.\,f.). \emph{[GA4] Data retention}. Analytics Help. \url{https://support.google.com/analytics/answer/7667196}
  \item Google. (s.\,f.). \emph{[GA4] BigQuery Export}. Analytics Help. \url{https://support.google.com/analytics/answer/9358801}
  \item Google. (s.\,f.). \emph{Conversions vs. key events in Google Analytics}. Analytics Help. \url{https://support.google.com/analytics/answer/13965727}
  \item Google. (s.\,f.). \emph{[GA4] Cookie usage on websites}. Analytics Help. \url{https://support.google.com/analytics/answer/11397207}
  \item Google. (s.\,f.). \emph{Consent mode reference}. Analytics Help. \url{https://support.google.com/analytics/answer/13802165}
  \item Google. (s.\,f.). \emph{Create a report} (Google Analytics Data API). Google for Developers. \url{https://developers.google.com/analytics/devguides/reporting/data/v1/basics}
  \item Google. (s.\,f.). \emph{Google Analytics demo account}. \url{https://goo.gle/ga-demo}
  \item Barsekh-Onji, A. (2026). \emph{Comportamiento del consumidor digital: datos, perfiles y modelos} (Caps.~1 a 4, material de curso). IPN, ESCA Unidad Santo Tomás.
\end{itemize}

\end{document}
